\section{Scalar comparison and its refinements}\label{rev:app-comparison}

Throughout this appendix, $E/k$ is an algebraic elliptic curve,
$k=\Qbar$, and the period matrix has determinant $2\pi i$.
Unless stated otherwise we assume the nonzero scalar relation
\[
 \omega(a\omega+bh)=\alpha\pi,\qquad a,b,\alpha\in k,
 \quad\alpha\ne0.
\]
The main text identifies its missing comparison coordinate.  We record
here the exact restricted hypothesis, useful analytic refinements,
and the limits of weaker replacement conditions.

\subsection{The restricted directions and known independence}
For the standard real branch put $t=\sqrt{1-z}$ and
$(e_1,e_2,e_3,e_4)=(0,1/4,1/2,0)$.  The Gauss--Manin coefficients
give
\[
 \left\{\frac ab:A/B\in\Q\right\}=t\Q-e_\ell
 \qquad(B\ne0).
\]
Indeed, substituting the explicit logarithmic derivatives of the
discriminants into $a/b=(A/B+u_s)/v_s$ yields a nonzero rational
multiple of $tA/B$, a rational multiple of $t$, and the constant
$-e_\ell$.  Consequently the precise restriction needed to exclude
non-CM rational slopes is transcendence of
\[
 \frac{\omega^2}{\pi},\qquad
 \frac{\omega h+(tq-e_\ell)\omega^2}{\pi}\quad(q\in\Q)
\]
at the specified non-CM fibers.  Multiplication by the nonzero
algebraic coefficient $b$ proves both directions of this equivalence.
This is weaker than QPI, but is an equivalent reformulation of the
rational-slope exclusion itself.

Two unconditional independence theorems clarify this restriction.
Chudnovsky gives algebraic independence of $h/\omega$ and
$\pi/\omega$.  Under the scalar relation, write $r=h/\omega$; then
\[
 \frac\pi\omega=\frac{\omega(a+br)}\alpha,
 \qquad
 \omega=\frac{\alpha(\pi/\omega)}{a+br}.
\]
The denominator is nonzero because $\alpha\ne0$.  Hence
$k(r,\pi/\omega)=k(r,\omega)=k(\omega,h)$, so $\omega,h$ would
be algebraically independent even under a hypothetical non-CM identity.
The theorem therefore does not rule out that identity.

Nesterenko's theorem, in \cite[Theorem~1.1]{Fonseca}, states
\[
 \operatorname{trdeg}_{\Q}\Q(q,E_2(\tau),E_4(\tau),E_6(\tau))\ge3
 \quad(\tau\in\mathfrak H).
\]
Put $\xi=\omega/(2\pi i)$ and $\upsilon=h/(2\pi i)$.
The exact lattice formulas
\[
 E_2=12\xi\upsilon,\qquad E_4=12g_2\xi^4,\qquad
 E_6=-216g_3\xi^6
\]
place the modular quantities in the three-generator field
$k(q,\xi,\upsilon)$.  That field therefore has transcendence degree
exactly three.  The proposed identity becomes
\[
 4\pi(a\xi^2+b\xi\upsilon)+\alpha=0,
\]
which only puts $\pi$ in $k(\xi,\upsilon)$.  It does not contradict
the independence of $q,\xi,\upsilon$.

The comparison fields of the main text give a concise account of the
conjectural deficit.  A relation $aP+bQ=\alpha$, where
$P=\omega^2/\pi$ and $Q=\omega h/\pi$, implies
\[
 \begin{aligned}
 \operatorname{trdeg}_k k(\tau,j',j'')&\le2,\\
 \operatorname{trdeg}_k K_E&\le3,\\
 \operatorname{trdeg}_k k(2\pi i,\tau,q,E_2,E_4,E_6)&\le4.
 \end{aligned}
\]
Each inequality follows from \eqref{eq:hierarchy-jets-field}--
\eqref{eq:hierarchy-modular-field}, adding at most one transcendence
degree for each extra generator.  These are precisely the deficits
contradicted by the respective numerical conjectures.

\subsection{Functional transcendence and specialization}
The unconditional Ax--Schanuel consequence in
\cite[Section~4.1]{EterovicMSCD} does not provide the numerical
MSCD assertion used here.  In that reference let
$C_j=\operatorname{jcl}(\varnothing)$.  It is an algebraically closed
field containing $k$, and Lemma~4.6 gives
\[
 j(\tau)\in C_j\Longrightarrow\tau\in C_j
 \Longrightarrow j(\tau),j'(\tau),j''(\tau)\in C_j.
\]
Its Proposition~4.7 measures transcendence over a $j$-closed field.
At an algebraic $j$-value all the displayed entries lie in $C_j$;
both the transcendence degree and the associated dimension in that
bound are zero.  Thus this established functional theorem gives no
positive numerical lower bound at the points required here.
A relation holding at one point also supplies no relation among its
derivatives on a neighborhood.  Similarly, an abstract algebraic
conjugation cannot be commuted with an infinite analytic sum without
a separate argument.

\subsection{A real refinement of cycle propagation}\label{sec:real-propagation}
Suppose now that the model and the coefficients are real algebraic.
Choose a fundamental basis with $\omega>0$ and
$\tau=r+iy$, where $r\in\{0,1/2\}$ and $y>0$.

\begin{proposition}[Mixed-cycle reality]\label{prop:real-propagation}
Under the scalar relation, the following are equivalent:
\begin{enumerate}
\item $E$ has CM;
\item $\omega_{m,n}(a\omega_{m,n}+bh_{m,n})$ is real for some
integers $m,n$ with $n(m+nr)\ne0$.
\end{enumerate}
At CM every such value satisfies
\[
 \frac{\omega_{m,n}(a\omega_{m,n}+bh_{m,n})}{\pi}
       =\alpha|m+n\tau|^2\in k\cap\mathbf R.
\]
\end{proposition}
\begin{proof}
By \eqref{eq:cycle-propagation}, the imaginary part of the value
divided by $\pi$ is
\[
 2n(m+nr)(\alpha y+b).
\]
The reality hypothesis gives $y=-b/\alpha\in k$.  Thus $\tau$ is
algebraic and Schneider's theorem proves CM.

Conversely, at CM the independence of $\omega$ and $\pi$ excludes
$b=0$.  Proposition~\ref{prop:ksv-exact-completion} then gives
$b=-\alpha y$.  For $L=m+n\tau$ we have
$\overline L=L-2iny$, so
$\alpha L^2+2ibnL=\alpha L\overline L$.  This proves the formula
and the converse.
\end{proof}

The condition $n(m+nr)\ne0$ is essential.  Values on a real or purely
imaginary cycle are automatically real in
\eqref{eq:cycle-propagation}.  In particular, Schwarz reflection,
which makes values on conjugate cycles conjugate, supplies no
mixed-cycle reality condition.

The same limitation can be seen before forming the quadratic value.
Write $\omega_2=r\omega+iv$ and $h_2=rh+iu$, with $u,v$ real.
Legendre gives $\omega u-vh=2\pi$, and the scalar relation gives
\[
 \omega(u-2a\omega/\alpha)
       -h(v+2b\omega/\alpha)=0.
\]
Conjugation repeats this equation.  If $v+2b\omega/\alpha=0$, then
$v/\omega$ is algebraic and Schneider forces CM.  Otherwise a
purported algebraic linearization must first establish algebraicity
of that same period ratio.

For completeness, the modular form of cycle transport is equally
exact.  If $\gamma=\left(\begin{smallmatrix}u&v\\c&d\end{smallmatrix}\right)
\in\operatorname{SL}_2(\Z)$ and $\lambda=c\tau+d$, then, changing
the period basis while keeping the algebraic model fixed,
\[
 P(\gamma\tau)=\lambda^2P(\tau),\qquad
 Q(\gamma\tau)=\lambda^2Q(\tau)+2ic\lambda.
\]
The first equation follows from $\omega_{d,c}=\lambda\omega$;
the second follows from Legendre.  Hence the transported affine
value is $\alpha\lambda^2+2ibc\lambda$, agreeing with
\eqref{eq:cycle-propagation} and supplying no automatic propagation.

\subsection{The almost-holomorphic coordinate}
In the same real normalization put
\[
 \kappa=\frac h\omega+\frac{\pi}{y\omega^2}.
\]
Under the scalar relation, algebraicity of $\kappa$ forces CM.  Indeed,
Legendre's identity gives the exact linear relation
\[
 h_2-\kappa\omega_2-\left(r+\frac{ib}{\alpha}\right)h
      +\left(r\kappa-\frac{ia}{\alpha}\right)\omega=0.
\]
If $\kappa$ were algebraic and $E$ non-CM, this would contradict
the algebraic linear independence of its four period entries
\cite[Section~5]{Waldschmidt}; the coefficient of $h_2$ is one.
Conversely, CM makes $\kappa$ algebraic.  To see the normalization
directly, diagonalize a nonreal CM endomorphism on de Rham cohomology.
Its antiholomorphic eigenvector has the form $\Xi-\kappa_0\Omega$
with $\kappa_0\in k$.  Thus
$h_j-\kappa_0\omega_j=K\overline{\omega_j}$.  Taking the determinant
with $(\omega_1,\omega_2)$ yields
\[
 K=-\frac\pi{y\omega^2},\qquad
 \kappa_0=\frac h\omega+\frac\pi{y\omega^2}=\kappa.
\]
This also proves that $\kappa$ is real.

For $g_2g_3\ne0$, the modular interpretation is
\[
 E_2^*=E_2-\frac3{\pi y},\qquad
 \frac{E_2^*E_4}{E_6}=-\frac{2g_2}{3g_3}\kappa.
\]
This follows by substitution of the lattice formulas.
Masser's theorem gives algebraicity of this almost-holomorphic value
at CM points \cite[Section~5]{Spence}.  Under the scalar relation
alone, however, one only has
\[
 \kappa=-\frac ab+\left(\frac\alpha b+\frac1y\right)
                      \frac\pi{\omega^2}\quad(b\ne0),
\]
whose coefficient $y$ is uncontrolled.  The exceptional cases
$g_2g_3=0$ already have CM.

\subsection{Coordinates of tensor completion}
Put $D=2\pi i$ and $k_0=\alpha/(2i)$.  The symmetric bilinear
coordinates of $a\Omega^2+b\Omega\Xi$, divided by its determinant
period, are
\[
 s_{ij}=\frac{a\omega_i\omega_j+
           (b/2)(\omega_i h_j+h_i\omega_j)}{D}.
\]
Legendre's identity gives
\begin{equation}\label{eq:comparison-coordinates}
 (s_{11},s_{12},s_{22})
       =(k_0,k_0\tau+b/2,k_0\tau^2+b\tau),\qquad
 s_{12}^2-s_{11}s_{22}=b^2/4.
\end{equation}
Thus algebraicity of $s_{12}$ is equivalent to algebraicity of
$\tau$, to CM, and to algebraicity of all three coordinates.  The
determinant relation is automatic and adds no missing coordinate.

When $b\ne0$, the tensor $t$ of
Proposition~\ref{prop:ksv-exact-completion} has, in the polynomial
Betti basis $(e_1^2,e_1e_2,e_2^2)/(e_1\wedge e_2)$, coordinates
\begin{equation}\label{eq:ksv-three-coordinates}
 t_{\mathrm B}
   =\left(\frac1c,\frac{2\tau+c}{c},
                       \frac{\tau(\tau+c)}c\right).
\end{equation}
The middle polynomial coefficient is twice the symmetric bilinear
coefficient.  The canonical adjoint identification sends $(x,y,z)$ to
\[
 \begin{pmatrix}-y/2&x\\-z&y/2\end{pmatrix}
       =\Pi_{\mathrm B}-\tfrac12I.
\]
This verifies the conventions in the intrinsic comparison square.

At CM the rational endomorphism algebra is the imaginary quadratic
field $K$, which has no nontrivial idempotent.  The complete projector
belongs to $K\otimes k\simeq k\times k$, and already exists with
coefficients in $K$; it is a $K$-linear combination of algebraic
endomorphisms.  Its entry field is exactly $K$ in a rational Betti
basis.  There is no assertion that it is a rational endomorphism.

\subsection{The local annihilator condition}
Set $T=\operatorname{Sym}^2H\otimes(\det H)^{-1}$,
$N=T\oplus\mathbf1$, and $\xi=(t,c^{-1})\in N_{\mathrm{dR}}$.
Let $\beta$ be the algebraic Betti functional which subtracts the
$\mathbf1$-coordinate from the first coordinate of $T$.  Equation
\eqref{eq:ksv-three-coordinates} gives
$\beta(\operatorname{comp}\xi)=0$.

\begin{proposition}[Local annihilator lift]\label{prop:ksv-local-lift}
Under the scalar identity with $b\ne0$, CM is equivalent to the
existence of a de Rham--Betti subobject $N'\subseteq N$ such that
\[
 \xi\in N'_{\mathrm{dR}},\qquad N'_{\mathrm B}\subseteq\ker\beta.
\]
\end{proposition}
\begin{proof}
For non-CM $E$, Theorem~\ref{thm:ksv-specialized} identifies the
algebraic-coefficient group with $\operatorname{GL}_2$.  Its adjoint
representation $T$ is irreducible and nontrivial.  Semisimplicity and
$\operatorname{Hom}(T,\mathbf1)=\operatorname{Hom}(\mathbf1,T)=0$
show that a subrepresentation of $T\oplus\mathbf1$ containing a
vector with nonzero components in both summands is the full object.
Both components of $\xi$ are nonzero.  Thus $N'=N$, contrary to
$N'_{\mathrm B}\subseteq\ker\beta$.
At CM, completion makes all coordinates of $\operatorname{comp}\xi$
algebraic.  The line spanned by $\xi$ and its comparison is a
de Rham--Betti subobject whose Betti line lies in $\ker\beta$.
\end{proof}

Proposition~3.6 of \cite{KSV} characterizes equality between the
numerical Zariski closure of a comparison point and its smallest
algebraic torsor by such annihilator lifting for every object in the
generated tensor category.  Proposition~5.8 proves lifting in the
abelian category generated by first cohomology of abelian varieties.
The adjoint tensor $T$ is not a first-cohomology object.  Theorem~5.13
identifies the de Rham--Betti torsor with the motivic torsor, without
asserting that the numerical Zariski closure equals that torsor.
The local lift above is consequently an exact additional premise.

This distinction also persists for the Kummer surface
$S=\operatorname{Km}(E\times E)$.  Its Picard rank is $19$ in the
non-CM case and $20$ in the CM case: the ranks $3$ and $4$ of
$E\times E$ acquire the sixteen exceptional divisors.  Corollary~6.17
of \cite{KSV} gives, respectively,
\[
 \Omega_S=\Omega_S^{\mathrm{mot}},\qquad
 Z_S=\Omega_S^{\mathrm{mot}}.
\]
Here $Z_S$ is the numerical Zariski closure and $\Omega_S$ is the
smallest algebraic subtorsor containing the comparison point.  The
rank-$19$ assertion does not identify $Z_S$ with that torsor.

\subsection{Formal obstructions to automatic completion}
For fixed $c\in k^\times$, every trace-one, rank-one projector with
upper-right entry $1/c$ is
\[
 \begin{pmatrix}x&1/c\\cx(1-x)&1-x\end{pmatrix},\qquad x\in\mathbb C.
\]
Indeed trace one fixes the second diagonal entry and determinant zero
fixes the lower-left entry; the Cayley--Hamilton identity then proves
idempotence.  Thus these algebraic conditions leave one free coordinate.
For example $c=-2i$, $\tau=iy$, $y>0$, gives
\[
 \begin{pmatrix}y/2&i/2\\i(y^2-2y)/2&1-y/2\end{pmatrix}.
\]
Conjugation of its entries followed by conjugation by
$\operatorname{diag}(1,-1)$ fixes the matrix, for every $y>0$.
Real structure and positivity therefore leave the same freedom;
actual CM completion would require $y=1$ by
\eqref{eq:ksv-cm-field}.  These are formal comparison objects, not
counterexamples over algebraic elliptic curves.

Connectedness of a torsor is likewise insufficient for genericity:
evaluation at $2$ on the coordinate ring $k[t,t^{-1}]$ of the connected
torsor $\mathbf G_m$ has kernel $(t-2)$.  More directly, set
\[
 f(M)=aw^2+bwh-\kappa\det M,\qquad \kappa\ne0,
 \quad M=\begin{pmatrix}w&h\\v&k_2\end{pmatrix}.
\]
At a zero of $f$ in $\operatorname{GL}_2$,
$f(\operatorname{diag}(1,r)M)=\kappa(1-r)\det M$.
Thus a numerical zero does not make the relation an identity on the
torsor.

\subsection{Why logarithm methods need additional data}
The analytic subgroup theorem applies to a logarithm with algebraic
exponential lying in an algebraic linear subspace of a commutative
algebraic Lie algebra.  The universal vector extension of $E$,
together with $\mathbf G_m$, provides the period logarithm
$(\omega,h,2\pi i)$ with exponential the identity.  Our relation
instead places it on
\[
 q(w,v,t)=aw^2+bwv-\frac\alpha{2i}t=0.
\]
The tangent hyperplane there is not algebraic.  Its normal vector is
$(2a\omega+bh,b\omega,-\alpha/(2i))$; since its last entry is nonzero
algebraic, an algebraic tangent hyperplane would make the other two
entries algebraic.  If $b\ne0$ this makes $\omega$ algebraic; if
$b=0$, then $a\ne0$ and the first entry has the same consequence.
Both contradict period transcendence
\cite[Corollary~3]{Waldschmidt}.

A bilinear central extension does not remove the obstruction.
Let $B:V\times V\to W$ be bilinear in characteristic zero.  The
associative law
\[
 (v,t)*(w,s)=(v+w,t+s+B(v,w))
\]
is commutative exactly when $B$ is symmetric, since its commutator is
$(0,B(v,w)-B(w,v))$.  In that case
$(v,t)\mapsto(v,t-B(v,v)/2)$ is an algebraic isomorphism to the additive
group.  Its differential at zero is the identity; hence the original
exponential is $(v,t)\mapsto(v,t+B(v,v)/2)$.  Algebraicity of its
endpoint still requires algebraicity of $v$, which fails for the
period vector.  This calculation is about global bilinear cocycles
on vector groups; it neither splits an elliptic universal extension
nor classifies general biextensions.

Nor can one repair this construction by reducing a quadratic coordinate
modulo a fixed discrete group.  Suppose a lattice $\Lambda$ spans a
complex vector space $L$, and a homogeneous quadratic form $q$ obeys
\[
 q(z+\lambda)-q(z)\in\Gamma
 \quad(z\in L,\ \lambda\in\Lambda)
\]
for a discrete additive subgroup $\Gamma\subset\mathbb C$.
Writing $q(z)=B(z,z)$, the affine function
$2B(z,\lambda)+B(\lambda,\lambda)$ has discrete image, so its linear
part is zero.  Since $\Lambda$ spans $L$, $B=0$.  The two elliptic
period--quasi-period vectors span their complex Lie algebra by
Legendre's determinant, so this obstruction applies to them.

Finally, the tensor itself lies outside first-cohomology logarithm
theory.  The Hodge types of $T=\operatorname{Sym}^2H^1(E)(1)$ are
$(1,-1),(0,0),(-1,1)$.  Cohomological realizations of 1-motives have
types among $(0,0),(1,0),(0,1),(1,1)$; their even-weight graded
pieces, and those of finite extensions of their Tate twists, are Tate.
Strictness of the Hodge and weight filtrations excludes every Tate
twist of the full $T$ as a subquotient of such an object.  For non-CM
$E$ there is not even a nonzero Hodge morphism between $T$ and these
objects.  Indeed
\[
 \operatorname{Hom}_{\mathrm{HS}}(\Q(0),T)
   =\operatorname{End}_{\mathrm{HS}}(H^1(E))_0
   =(\operatorname{End}(E)\otimes\Q)_0=0.
\]
The middle identification follows by clearing denominators in a Hodge
map and interpreting it as a holomorphic homomorphism of complex tori.
Duality excludes Tate quotients as well.  Any morphism to or from the
stated mixed objects would factor through such a Tate piece by
strictness, proving the assertion.

Rigidified $\mathbf G_m$-biextensions of $E\times E$ are classified
by $\operatorname{Hom}(E,E^\vee)$: restrict to $\{x\}\times E$ to
obtain a point of $E^\vee$; the biextension law makes this map a
homomorphism, and pullback of the Poincar\'e bundle is inverse to this
construction.  Under a principal polarization these are endomorphisms
of $E$.  In the non-CM case their rational first Chern tensors are
multiples of the polarization and have zero projection to $T$.
A nonzero symmetric-square Chern tensor would therefore already force
CM.  It would be algebraic by the Lefschetz $(1,1)$ theorem.  Absolute
Hodge theory starts with such a complete Hodge class; it does not
construct the class from one scalar period value.  Logarithms of
algebraic points on a biextension require further data, whose
algebraicity must be established before a logarithm theorem applies.
